\documentclass[11pt,a4paper]{article}
\usepackage[margin=2.5cm]{geometry}
\usepackage{amsmath,amssymb,amsthm}

\newcommand{\R}{\mathbb{R}}
\newcommand{\E}{\mathbb{E}}
\newcommand{\norm}[1]{\left\| #1 \right\|}
\newcommand{\ip}[2]{\left\langle #1, #2 \right\rangle}
\newcommand{\xs}{x^\star}

\title{Homework 00.11.26 -- Exercises 31--37}
\author{Gor Melkumyan}
\date{}

\begin{document}
\maketitle

\section*{Exercise 31 (Decomposition of Bregman divergence)}

Let $f = \frac{1}{n}\sum_{i=1}^n f_i$. Prove that
\[
  D_f(x,y) = \frac{1}{n}\sum_{i=1}^n D_{f_i}(x,y).
\]

\begin{proof}
We start from the definition of the Bregman divergence of $f$ and substitute
$f = \frac{1}{n}\sum_{i=1}^n f_i$. Since the gradient is linear, we also have
$\nabla f(y) = \frac{1}{n}\sum_{i=1}^n \nabla f_i(y)$, so
\[
  D_f(x,y)
  = f(x) - f(y) - \ip{\nabla f(y)}{x-y}
  = \frac{1}{n}\sum_{i=1}^n f_i(x) - \frac{1}{n}\sum_{i=1}^n f_i(y)
    - \ip{\frac{1}{n}\sum_{i=1}^n \nabla f_i(y)}{x-y} .
\]
The inner product is linear in its first argument, so the last term is also an average over $i$.
Collecting the terms with the same index $i$ under one sum, we see that each bracket is exactly
the Bregman divergence of $f_i$:
\[
  D_f(x,y)
  = \frac{1}{n}\sum_{i=1}^n \underbrace{\bigl(f_i(x) - f_i(y) - \ip{\nabla f_i(y)}{x-y}\bigr)}_{D_{f_i}(x,y)}
  = \frac{1}{n}\sum_{i=1}^n D_{f_i}(x,y). \qedhere
\]
\end{proof}

\section*{Exercise 32 (Theorem 38)}

Assume that
\begin{enumerate}
  \item $f$ is $\mu$-convex.
  \item Each $f_i$ is convex and $L_i$-smooth.
\end{enumerate}
Then SGD-NS with stepsize $0 < \gamma \le \frac{1}{2A''}$, where
$A'' = \max_i \frac{L_i}{np_i}$, satisfies
\[
  \E\bigl[\norm{x^k - \xs}^2\bigr] \le (1-\gamma\mu)^k\norm{x^0 - \xs}^2 + \frac{2\gamma\sigma_\star^2}{\mu},
\]
where
\[
  \sigma_\star^2 = \mathrm{Var}\bigl[g(\xs)\bigr]
  = \E\bigl[\norm{g(\xs)}^2\bigr] - \norm{\E\bigl[g(\xs)\bigr]}^2
  = \left(\frac{1}{n^2}\sum_{i=1}^n \frac{\norm{\nabla f_i(\xs)}^2}{p_i}\right) - \norm{\nabla f(\xs)}^2 .
\]

\begin{proof}
The theorem is a consequence of Theorem 33 and Lemma 37, so the plan is to check that the
assumptions of Theorem 33 hold for SGD-NS and to compute the constants $A$ and $C$.
The function $f$ is $\mu$-convex by assumption, and the stochastic gradient
$g(x) = \frac{\nabla f_i(x)}{np_i}$ of SGD-NS is unbiased, since
$\E\bigl[g(x)\bigr] = \sum_{i=1}^n p_i \frac{\nabla f_i(x)}{np_i} = \nabla f(x)$.
So it remains to check the AC assumption.

From Lemma 37, using the convexity and $L_i$-smoothness of $f_i$, we have that
\[
  \E\bigl[\norm{g(x) - g(y)}^2\bigr] \le 2A''\,D_f(x,y), \qquad \text{where } A'' = \max_i \frac{L_i}{np_i}.
\]
Moreover, for $G(x,y) = \E\bigl[\norm{g(x) - \nabla f(y)}^2\bigr]$ we have
\[
  G(x,y) \le 4A''\,D_f(x,y) + 2\,\mathrm{Var}\bigl[g(y)\bigr] \qquad \forall x, y \in \R^d ,
\]
which for $x = x^k$ and $y = \xs$ is the AC assumption.

So, substituting everything into Theorem 33, we get
\[
  \E\bigl[\norm{x^k - \xs}^2\bigr] \le (1-\gamma\mu)^k\norm{x^0 - \xs}^2 + \frac{\gamma C}{\mu},
\]
where $0 < \gamma \le \frac{1}{A} = \frac{1}{2A''}$ with
\[
  A = 2A'' = 2\max_i \frac{L_i}{np_i} \qquad \text{and} \qquad C = 2\,\mathrm{Var}\bigl[g(\xs)\bigr] = 2\sigma_\star^2 . \qedhere
\]
\end{proof}

\section*{Exercise 33}

Prove that
\[
  \sigma_\star^2 = \E\bigl[\norm{g(\xs) - \nabla f(\xs)}^2\bigr] = 0
  \quad\text{iff}\quad \nabla f_i(\xs) = np_i\nabla f(\xs) \quad \forall i = 1,\dots,n.
\]

\begin{proof}
We use the formula for $\sigma_\star^2$ from Exercise 32,
\[
  \sigma_\star^2 = \left(\frac{1}{n^2}\sum_{i=1}^n \frac{\norm{\nabla f_i(\xs)}^2}{p_i}\right) - \norm{\nabla f(\xs)}^2 ,
\]
and prove the two directions separately.

\medskip
\noindent\textbf{1) ($\Leftarrow$)}
Assume that $\nabla f_i(\xs) = np_i\nabla f(\xs)$ for all $i$. Substituting this into the formula,
we get $\norm{np_i\nabla f(\xs)}^2/p_i = n^2p_i\norm{\nabla f(\xs)}^2$, so the factor $\frac{1}{n^2}$
cancels, and since the $p_i$ sum up to one, the two terms are equal:
\begin{align*}
  \left(\frac{1}{n^2}\sum_{i=1}^n \frac{\norm{\nabla f_i(\xs)}^2}{p_i}\right) - \norm{\nabla f(\xs)}^2
  &= \left(\frac{1}{n^2}\sum_{i=1}^n \frac{\norm{np_i\nabla f(\xs)}^2}{p_i}\right) - \norm{\nabla f(\xs)}^2 \\
  &= \Biggl(\underbrace{\sum_{i=1}^n p_i}_{1}\Biggr)\norm{\nabla f(\xs)}^2 - \norm{\nabla f(\xs)}^2 = 0
  \;\Rightarrow\; \sigma_\star^2 = 0 .
\end{align*}

\medskip
\noindent\textbf{2) ($\Rightarrow$)}
Let $\sigma_\star^2 = 0$, and let's assume that there are $p_i' > 0$ such that $\sum_i p_i' = 1$ and
$\nabla f_i(\xs) = np_i'\nabla f(\xs)$. If $\nabla f(\xs) = 0$, then $\nabla f_i(\xs) = 0 = np_i\nabla f(\xs)$
and there is nothing to prove, so let $\nabla f(\xs) \neq 0$. Substituting $\nabla f_i(\xs)$ into the
formula and taking $\norm{\nabla f(\xs)}^2$ out of the brackets, we get
\[
  \sigma_\star^2 = \norm{\nabla f(\xs)}^2\left(\frac{1}{n^2}\sum_{i=1}^n \frac{(p_i'n)^2}{p_i} - 1\right) = 0
  \iff \sum_{i=1}^n \frac{p_i'^2}{p_i} = 1 .
\]
One such solution for $p_i'$ is $p_i' = p_i$. It also turns out that it is the only solution:
indeed, using Exercise 38 with $a_i := p_i'^2$ and $p_i := p_i$, we get the result.
\end{proof}

\section*{Exercise 34}

Prove part (iii) of Lemma 44.

\paragraph{Lemma 44.}
Let $f_1,\dots,f_n : \R^d \to \R$ be differentiable and $f(x) := \frac{1}{n}\sum_{i=1}^n f_i(x)$.
\begin{enumerate}
  \item[(i)] If $f_i$ are $L_i$-smooth, then $f$ is $L$-smooth with $L \le \frac{1}{n}\sum_{i=1}^n L_i$.
  \item[(ii)] If $f$ is $L$-smooth and $f_i$ are convex, then $f_i$ is $L_i$-smooth with $L_i \le nL$.
  \item[(iii)] If $f$ is $L'$-smooth and $f_i$ are convex and $L_i'$-smooth, then there exist
    $L \le L'$ and $L_i \le L_i'$ $\forall i$ such that $f$ is $L$-smooth, $f_i$ is $L_i$-smooth
    $\forall i$, and $L \le \max_i L_i \le nL$.
\end{enumerate}

\begin{proof}[Proof of (iii)]
The idea is to take the smallest smoothness constants of $f$ and $f_i$ and then to apply
parts (i) and (ii) to them. Let's take
\[
  L = \min\Bigl(L',\ \inf\bigl\{L^* \in \R^+ : f \text{ is } L^*\text{-smooth}\bigr\}\Bigr)
\]
and
\[
  L_i = \min\Bigl(L_i',\ \inf\bigl\{L_i^* \in \R^+ : f_i \text{ is } L_i^*\text{-smooth}\bigr\}\Bigr).
\]
Now, using (i) for the $L_i$-smooth functions $f_i$, we get that $f$ is $L''$-smooth and
$L'' \le \frac{1}{n}\sum_i L_i$, but because $L \le L''$ (recall that $L$ is the smallest smoothness
constant of $f$), we get that
\[
  L \le \frac{1}{n}\sum_{i=1}^n L_i \le \max_i L_i .
\]
Now, using (ii) for the $L$-smooth function $f$ and the convex functions $f_i$, we get that $f_i$ is
$L_i''$-smooth and $L_i'' \le nL$, but because $L_i \le L_i''$ (again, $L_i$ is the smallest smoothness
constant of $f_i$), we get that $L_i \le nL$, so $\max_i L_i \le nL$, which proves Lemma 44 (iii):
\[
  L \le \frac{1}{n}\sum_{i=1}^n L_i \le \max_i L_i \qquad \text{and} \qquad \max_i L_i \le nL . \qedhere
\]
\end{proof}

\section*{Exercise 35}

Give an example of $L_i$-smooth functions $f_i$ for which
$\frac{1}{n}\sum_i L_i \ll \max_i L_i$.

\begin{proof}[Solution]
The idea is to make one function ``much less smooth'' than all the others: one large $L_1$ is enough to
make $\max_i L_i$ large, while in the average it is divided by $n$. Let $n = 1000$ and
\[
  f_1 = 500x^2, \qquad f_2 = x^2, \qquad f_i = \tfrac{1}{2}x^2 \quad \forall i = 3,\dots,n.
\]
A function $cx^2$ with $c > 0$ has the derivative $2cx$, which is $2c$-Lipschitz, so it is $2c$-smooth.
Hence $L_1 = 1000$, $L_2 = 2$ and $L_i = 1$ for $i \ge 3$. In this case
\[
  \frac{1}{n}\sum_{i=1}^n L_i = \frac{1}{1000}\cdot\Bigl(1000 + 2 + \sum_{i=3}^{1000} 1\Bigr)
  = \frac{2000}{1000} = 2
\]
and $\max_i L_i = 1000$, so we have
\[
  \max_i L_i = 1000 \gg 2 = \frac{1}{n}\sum_{i=1}^n L_i . \qedhere
\]
\end{proof}

\section*{Exercise 36}

Let $a_1,\dots,a_n > 0$. Prove that the solution to the optimization problem
\[
  D_{\Delta_n} := \min_{p \in \Delta_n} \max_i \frac{a_i}{p_i}
\]
is given by $p_i = \frac{a_i}{\sum_j a_j}$ for all $i = 1,\dots,n$.

\begin{proof}
The proof consists of two steps: first we compute the value of the objective at $p^*$ (here it is a
partition that is defined in the exercise), and then we show that any other $p' \in \Delta_n$ gives a
strictly larger value. Note that $p^* \in \Delta_n$, since all $p_i^* > 0$ and $\sum_i p_i^* = 1$.

First of all, note that for $p_i^* = \frac{a_i}{\sum_j a_j}$ we have
\[
  D := \max_i \frac{a_i}{p_i^*} = \max_i \Bigl(\sum_j a_j\Bigr) = \sum_j a_j ,
\]
i.e.\ all the ratios $\frac{a_i}{p_i^*}$ are equal to the same number $\sum_j a_j$.

Let's take another $p' \in \Delta_n$ and calculate $D' = \max_i \frac{a_i}{p_i'}$. If $p' \ne p^*$,
then there is an $i$ with $p_i' \ne p_i^*$. If $p_i' < p_i^*$, then let's assign $k = i$, and if
$p_i' > p_i^*$, then, since $\sum_j p_j' = \sum_j p_j^* = 1$, there is $j \ne i$ with $p_j' < p_j^*$, and
we will assign $k = j$. So in both cases there is a $k$ with $p_k' < p_k^*$. Since $a_k > 0$ and
$0 < p_k' < p_k^*$, decreasing the denominator increases the ratio, and we get
\[
  D = \sum_j a_j = \frac{a_k}{p_k^*} < \frac{a_k}{p_k'} \le \max_i \frac{a_i}{p_i'} = D' .
\]
So for any $p' \ne p^*$ we have $D < D'$, which means that $p^*$ is the (unique) solution.
\end{proof}

\section*{Exercise 37 (Importance sampling 2 -- simple version)}

Let $a, b > 0$. Prove that the solution to the optimization problem
\[
  \min_{0 < p < 1} \frac{a}{p} + \frac{b}{1-p}
\]
is given by $p = \frac{\sqrt{a}}{\sqrt{a} + \sqrt{b}}$.

\begin{proof}
We find the minimiser by looking at the first and second derivatives of the objective.
Let $\varphi(p) := \frac{a}{p} + \frac{b}{1-p}$ for $p \in (0,1)$. Then
\[
  \varphi'(p) = -\frac{a}{p^2} + \frac{b}{(1-p)^2}, \qquad
  \varphi''(p) = \frac{2a}{p^3} + \frac{2b}{(1-p)^3} > 0 ,
\]
so $\varphi$ is strictly convex on $(0,1)$, and a point where $\varphi'$ equals $0$ is its unique
minimiser. Equating the derivative to $0$ and using $p > 0$, $1-p > 0$ to take square roots, we get
\[
  \varphi'(p) = 0
  \iff \frac{a}{p^2} = \frac{b}{(1-p)^2}
  \iff \frac{\sqrt{a}}{p} = \frac{\sqrt{b}}{1-p}
  \iff \sqrt{a}\,(1-p) = \sqrt{b}\,p
  \iff p = \frac{\sqrt{a}}{\sqrt{a} + \sqrt{b}} .
\]
Since $a, b > 0$, this $p$ lies in $(0,1)$, so it is the solution. Using
$1 - p = \frac{\sqrt{b}}{\sqrt{a} + \sqrt{b}}$, the optimal value is
\[
  \varphi\left(\frac{\sqrt{a}}{\sqrt{a} + \sqrt{b}}\right)
  = \sqrt{a}\,\bigl(\sqrt{a} + \sqrt{b}\bigr) + \sqrt{b}\,\bigl(\sqrt{a} + \sqrt{b}\bigr)
  = \bigl(\sqrt{a} + \sqrt{b}\bigr)^2 . \qedhere
\]
\end{proof}

\end{document}
